\chapter{Los datos: de dónde sale cada número}
\label{cap:s-datos}

\section{La regla de oro: no inventar}
Todo sale de \textbf{fuentes oficiales o abiertas}, descargadas \textbf{con programas} (no a mano) y guardadas sin
tocar. Si un dato no existe, \textbf{no se inventa, no se rellena con un promedio y no se ``simula para probar''}: se dice
``sin dato'' y se explica por qué.

\begin{porque}
Descargar a mano desde una página (copiar y pegar, o ``exportar a Excel'') es más rápido la primera vez, pero no se puede
repetir ni comprobar. Con un programa, cualquiera puede volver a descargar todo con un comando y obtener lo mismo. Además,
el programa anota de dónde vino cada archivo, cuándo y cuántos registros tiene.
\end{porque}

También hay una regla de lo que \textbf{no} se usa: nada copiado de Google Maps ni de TripAdvisor, porque sus términos de
uso prohíben extraer sus datos con programas.

\section{Las fuentes, una por una}
Se descargaron \textbf{353 archivos con 8,134,802 registros} (824 MB). Estas son las fuentes, de la más grande a la más
chica:

\begin{center}
\small
\begin{tabularx}{\linewidth}{>{\raggedright\arraybackslash}p{0.27\linewidth}rY}
\encabezado \h{Fuente (quién la publica)} & \h{Registros} & \h{Qué es y para qué se usó} \\
DENUE (INEGI) & 6,138,075 & Todos los negocios del país, con su giro y su ubicación. Para saber cuántos hoteles,
restaurantes y museos hay en cada lugar, y qué recomendar en la página. \\
\filagris Clima (Open-Meteo) & 767,016 & El clima de cada día desde 1950 y de cada hora desde 2019, en 8 puntos. Para ver
si la lluvia aleja visitantes y detectar semanas de clima raro. \\
Nacionalidades (SECTUR) & 521,364 & Cuántos extranjeros llegan en avión a cada aeropuerto y de qué país. Para conocer a
la viajera extranjera. \\
\filagris Compendio 2024 (SECTUR) & 295,191 & Estadísticas generales del turismo; de aquí salen las estrellas oficiales de
los hoteles. \\
Rest-Mex 2025 (concurso académico) & 208,051 & Opiniones de turistas con su calificación de 1 a 5. Para saber qué
molesta y qué enamora. \\
\filagris INAH (SECTUR) & 71,562 & Visitantes por mes a cada zona arqueológica y museo. La serie más importante del sur. \\
Huracanes HURDAT2 (NOAA) & 57,512 & La trayectoria de todas las tormentas del Atlántico desde 1851. \\
\filagris Ocupación hotelera (SECTUR) & 34,164 & La ocupación de cada semana de 7 centros del norte. La única ocupación
medida en 2025–2026. \\
Vuelos AFAC & 19,578 & Pasajeros nacionales por aerolínea. \\
\filagris Tipo de cambio (FRED) & 9,531 & Pesos por dólar cada día, y la inflación de Estados Unidos. \\
SITUR-Q (gobierno del estado) & 6,607 & Tren Maya, cruces desde Belice, cruceros, cuartos de hotel y ocupación hasta 2024. \\
\filagris Cruceros (SECTUR) & 3,897 & Pasajeros de crucero por puerto. \\
Censo 2020 (INEGI) & 2,243 & La población de cada pueblo de Quintana Roo. \\
\filagris Uso de internet ENDUTIH (INEGI) & PDF & Cuánta gente usa redes y mensajería. \\
Mapa del estado (GeoJSON) & 11 & Los 11 municipios, para comprobar que cada lugar y cada foto están donde dicen. \\
\bottomrule
\end{tabularx}
\normalsize
\end{center}

Además: los \textbf{costos de anuncios} de Google y Facebook (WordStream y LocaliQ, guardados con captura de pantalla
como evidencia) y las \textbf{fotos} de Wikimedia Commons, con licencia libre y crédito a su autor.

\begin{error}
\textbf{El conteo del DENUE estaba inflado.} La primera cuenta decía 6,139,989 negocios. Al limpiarlos, Spark contó
6,138,075, y un segundo programa independiente le dio la razón. La diferencia era de 1,914: cada uno de los 33 archivos
comprimidos del DENUE trae, además de los negocios, un diccionario de datos de 58 renglones, y el conteo los sumaba como
si fueran negocios ($33\times58=1{,}914$). Se corrigió el programa y los documentos. Con el Censo pasó lo mismo: decía
2,257 pueblos y son 2,243.
\end{error}

\section{Los huecos: lo que no existe}
El hallazgo más importante de la descarga fue un hueco: \textbf{el gobierno del estado dejó de publicar la ocupación
hotelera en 2025, para todos los destinos}, incluso Cancún. Para 2025 y 2026 la única ocupación medida es la semanal de
SECTUR, y solo para 7 centros del norte.

\textbf{Consecuencia}: la ocupación del sur en 2025–2026 se muestra como ``sin dato oficial''. No se calcula con un
modelo, porque sería poner una cifra inventada donde nadie midió. En su lugar, el sur se mide con lo que \emph{sí} llega
hasta 2026: la gente que baja del Tren Maya, los cruces desde Belice y los visitantes a las zonas arqueológicas.

Otros huecos declarados:
\begin{itemize}
\item De qué estado del país viene el turista mexicano, su edad y su ingreso.
\item Cuánta gente compra un viaje después de ver un anuncio de turismo en Facebook.
\item Reseñas de los cinco lugares: Rest-Mex solo trae Tulum, Isla Mujeres y Bacalar.
\item Las tormentas de 2026 (la NOAA todavía no las publica).
\item Visitantes de Maya Ka'an y de la Laguna Milagros: ninguna fuente los cuenta.
\end{itemize}

\section{Cero no es lo mismo que ``sin dato''}
\begin{ejemplo}[ · los aeropuertos en 2025]
En los datos del gobierno, los cuatro aeropuertos del estado (Cancún, Cozumel, Chetumal y Tulum) dicen ``0 pasajeros''
todos los meses de 2025. Si se tomara como cero, la página diría ``nadie llegó en avión en 2025''. Pero ese mismo año la
ocupación semanal de Cancún promedió 72.2\,\%: es imposible que nadie haya volado. Lo que pasó es que \textbf{dejaron de
publicar}.

\textbf{Regla}: si en un mes \emph{todos} los aeropuertos marcan cero, ese mes es hueco. Si al menos uno trae dato, los
ceros se respetan (podría ser un aeropuerto cerrado de verdad). Con esta regla, 114 renglones pasaron a ser hueco, y la
página usa 2024 como último año completo del avión, y lo dice.
\end{ejemplo}

Lo mismo con otras cosas: una ocupación con 0 cuartos disponibles es imposible (hueco); en cambio, una zona arqueológica
con 0 visitantes \emph{sí} puede ser real, porque estaba cerrada, y se conserva marcada.

\section{Cuando dos fuentes no coinciden}
\begin{ejemplo}[ · Isla Mujeres]
Las noticias de 2026 decían que Isla Mujeres estaba al 93.5\,\% y al 95\,\% de ocupación. La fuente oficial (SECTUR)
registra \textbf{74.7\,\%} en febrero y \textbf{43.2\,\%} en abril. Además, SECTUR y el gobierno del estado no miden igual:
en Isla Mujeres hay 19.3 puntos de diferencia en promedio.
\end{ejemplo}
\textbf{Reglas que se fijaron}:
\begin{enumerate}
\item El dato oficial va antes que la nota de prensa.
\item Cada lugar usa \textbf{una sola fuente} en toda su historia, y la diferencia con la otra se declara. Así un cambio
  en el tiempo es real y no un salto entre fuentes.
\end{enumerate}
Con los cruceros pasa algo parecido: en Mahahual, el gobierno del estado cuenta siempre entre 11\,\% y 35\,\% más
cruceristas que SECTUR. Se usa una y se declara la diferencia.

\section{Privacidad}
No se guarda ningún dato personal. Del DENUE se quitan la razón social, el teléfono y el correo de cada negocio, porque
ningún modelo los necesita. Las reseñas de Rest-Mex no traen nombres, y las llegadas de extranjeros vienen sumadas por
país, no por persona.
